\documentclass[10pt]{article}
\usepackage[margin=0.88in]{geometry}
\usepackage[T1]{fontenc}
\usepackage{lmodern,amsmath,amssymb,amsthm,mathtools,microtype,booktabs}
\usepackage[hidelinks]{hyperref}
\newtheorem{theorem}{Theorem}
\newtheorem{lemma}{Lemma}
\newtheorem{corollary}{Corollary}
\newcommand{\E}{\mathbb E}
\newcommand{\Prb}{\mathbb P}
\newcommand{\Var}{\operatorname{Var}}
\newcommand{\Poi}{\operatorname{Poi}}
\newcommand{\eff}{\operatorname{eff}}
\newcommand{\ind}{\mathbf 1}
\DeclareMathOperator{\arcosh}{arcosh}
\DeclareMathOperator{\arsinh}{arsinh}
\DeclareMathOperator{\artanh}{artanh}
\title{Explicit Constants for Support-Size Testing:\\An Exact-Certificate Analysis}
\author{}
\date{}
\begin{document}
\maketitle
\begin{abstract}
We revisit both the sampling constant and the branch threshold in Wang's
support-size tester. Keeping all other scales, degrees, weights and decision
rules fixed, we establish the sufficient pair
$C_0=144(\log 2)^2=69.1852340042\ldots$, $L_0=16$.
Both the promised binary test and the joint support-bound guarantee succeed
with probability greater than $2/3$, uniformly over all finite or countably
supported distributions, all positive integers $n$, and all $0<\epsilon<1$.
The integer pair $(70,16)$ is also certified. The pair $(C_0,15)$ has success
greater than $3/5$. A practical alternative, $(70,20)$, has success at least
$0.92$, so five independent runs and majority/median aggregation have failure
below $0.005$ at total expected sampling coefficient $350$.
The proof combines individual coefficient bounds, a range-based variance
bound, one-sided concentration, and exact rational certificates on
$15\le L\le32$; the entire unbounded range $L\ge32$ is handled analytically.
The certificate verifier uses only integer and rational arithmetic, with
explicitly proved transcendental enclosures. We also give a bounded-sample
implementation with the same stated success guarantees. These are sufficient
constants, not globally minimal constants.
\end{abstract}

\section{Model, unchanged construction, and precise guarantees}
Write $C_*=144(\log2)^2$ throughout.
Let $p$ be a probability distribution on a finite or countable alphabet.
Write $S=|\operatorname{supp}(p)|$, possibly infinite, and let $p_{(i)}$ be
its nonincreasing rearrangement, padded with zeros when needed. Set
\[
 \Delta_r=\sum_{i>r}p_{(i)},\qquad
 \eff_\epsilon(p)=\min\{k\ge1:\Delta_k\le\epsilon\}.
\]
Here $\eff_\epsilon(p)$ is finite because $\Delta_k\downarrow0$.
All logarithms are natural unless explicitly subscripted. The separation
parameter $\epsilon$ is not the algorithm's failure probability.
The distributions in this paper are discrete; no claim about arbitrary
non-atomic probability spaces is needed.

\begin{lemma}[Distance interpretation]\label{lem:tv}
For an integer $n\ge1$, let $\mathcal P_n$ be the set of distributions supported on at most $n$
labels. With $d_{\rm TV}(p,q)=\tfrac12\sum_i|p_i-q_i|$,
$\inf_{q\in\mathcal P_n}d_{\rm TV}(p,q)=\Delta_n$.
\end{lemma}
\begin{proof}
For any support set $A$ of size at most $n$,
$d_{\rm TV}(p,q)\ge p(A^c)\ge\Delta_n$.
Conversely, if $S\le n$, choose $q=p$. If $S>n$, take $A$ to contain the
$n$ largest positive masses and normalize the
restriction of $p$ to $A$. Its mass is positive, and the total removed mass
and total added mass both equal $\Delta_n$, so its TV distance is $\Delta_n$.
\end{proof}

For an integer $n\ge1$ and $0<\epsilon<1$, define
\begin{equation}\label{eq:model}
 L=1+\log n,\qquad m=\frac{Cn}{\epsilon L},\qquad
 M\sim\Poi(m),\qquad N_i\sim\Poi(mp_i)\text{ independently},\qquad
 D=\sum_i\ind\{N_i>0\}.
\end{equation}
This is equivalent to $M$ independent draws from $p$ with an independent
Poisson sample count. A real Poisson mean is allowed; no rounding is needed.
For completeness, the joint probability generating function of any finite
collection of counts is
$\exp(m\sum_i p_i(z_i-1))$, which factors into the Poisson generating
functions. The total count is finite almost surely. The estimator below
therefore involves finitely many nonzero summands, even on an infinite
alphabet. Probabilities in all theorems are over these counts. We analyze
the ideal statistic with exact real arithmetic; bit complexity and
floating-point approximations are separate questions.
If $L<L_0$, let $\widetilde S=D$. If $L\ge L_0$, use
\begin{align}
 a_0&=\frac{\epsilon}{8n},& b&=\frac{64L}{m},&
 J&=\left\lfloor\log_2\frac{b}{2a_0}\right\rfloor,
 &a_j&=2^ja_0,\label{eq:parameters}\\
 \delta_j&=2^{-j-2},&q_j&=(j+3)\log2,&
 k_j&=\left\lceil\sqrt{b/a_j}\,q_j\right\rceil
 &&(0\le j\le J).\notag
\end{align}
For each scale, define
\begin{align}
 P_j(x)&=-\rho_j T_{k_j}\left(\frac{b+a_j-2x}{b-a_j}\right)
       =\sum_{\ell=0}^{k_j}c_{j,\ell}x^\ell,
 &\rho_j&=T_{k_j}\left(\frac{b+a_j}{b-a_j}\right)^{-1},\label{eq:weights}\\
 f_j(0)&=0,&f_j(\ell)&=1+c_{j,\ell}\ell!/m^\ell\quad(1\le\ell\le k_j),\notag\\
 f_j(\ell)&=1\quad(\ell>k_j),&Z_j&=\sum_i f_j(N_i),\notag\\
 \widetilde S&=\max\left\{D,\max_{0\le j\le J}
                    \frac{Z_j-\delta_j n}{1+\delta_j}\right\}.
 \label{eq:statistic}
\end{align}
Here $T_k(\cos\theta)=\cos(k\theta)$ is the Chebyshev polynomial
\cite[18.5.1]{dlmf}.
Return Yes exactly when $\widetilde S\le n$, and return the integer bound
\begin{equation}\label{eq:output}
 \widehat S=\min\{\max\{\lceil\widetilde S\rceil,1\},n\}.
\end{equation}
Only $C$ and $L_0$ differ from the construction of Wang~\cite[Section 3]{wang}.
The definitions at the branch boundary use $L\ge L_0$, and equality in the
test statistic is resolved as Yes. All degree and scale roundings are part
of the algorithm.
The binary problem has the two promises $S\le n$ and $\Delta_n\ge\epsilon$.
No binary correctness label is prescribed when $S>n$ and
$\Delta_n<\epsilon$; the support-bound guarantee below nevertheless applies
to those distributions as well.

\begin{theorem}[Certified sufficient pairs]\label{thm:main}
For every $p,n,\epsilon$ as above, every row in Table~\ref{tab:constants}
has the following three guarantees, with the indicated common upper bound $B$:
\begin{align}
 S\le n&\implies\Prb(\widetilde S>n)<B,\label{eq:null}\\
 \Delta_n\ge\epsilon&\implies\Prb(\widetilde S\le n)<B,\label{eq:alternative}\\
 \Prb\{\min(\eff_\epsilon(p),n)\le\widehat S\le S\}&>1-B.
 \label{eq:joint}
\end{align}
In particular, $(144(\log2)^2,16)$ and $(70,16)$ have success greater than
$2/3$ for both tasks. The respective pairs with $L_0=15$ have success
greater than $3/5$. Every displayed finite decimal $B$ is an exact rational
upper bound, proved by the integer-arithmetic certificate in
Appendix~\ref{app:arithmetic}. The table gives three individual guarantees;
simultaneous correctness of both outputs is addressed below.
\end{theorem}

\begin{table}[h]
\centering
\begin{tabular}{rrrl}
\toprule
$C$ & $L_0$ & rational failure upper bound $B$ & success guarantee\\
\midrule
$144(\log2)^2$ &15&0.392089&$3/5$\\
$144(\log2)^2$ &16&0.317480&$2/3$\\
$144(\log2)^2$ &20&0.087747&$0.91$\\
70&15&0.376264&$3/5$\\
70&16&0.282612&$2/3$\\
70&19&0.090185&$0.90$\\
70&20&0.072224&$0.92$\\
\bottomrule
\end{tabular}
\caption{Uniform sufficient bounds. Each probability statement includes all
distributions in its stated class, not just distributions used in experiments.}
\label{tab:constants}
\end{table}

\section{Expectation and weight bounds}
We provide the analytic ingredients in full. Their improvements concern
the analysis, not the estimator.

\begin{lemma}[Expectation]\label{lem:expect}
For a generic scale $0<a<b\le1$, let $P,\rho,f$ be as in
\eqref{eq:weights}. If $m\sqrt{b(b-a)}\ge k$, then
\[
 (1-\rho)\min\{x/a,1\}\le Q(x):=\E f(\Poi(mx))\le1+\rho
 \qquad(x\ge0).
\]
For the construction, $\rho_j\le\delta_j^2/2$.
\end{lemma}
\begin{proof}
The Poisson series yields $Q(x)=1+e^{-mx}P(x)$. On $[0,a]$, $P$ is
nonpositive, increasing, and concave: $T_k$ and its first two derivatives
are nonnegative on $[1,\infty)$ (expand about its roots, all in $(-1,1)$).
The chord from $P(0)=-1$ to $P(a)=-\rho$ proves the lower bound, since
$Q\ge1+P$ there; also $Q\le1$.
On $[a,b]$, $|P|\le\rho$. For $x\ge b$,
\[
 |P(x)|\le\rho\exp\left(2k\sqrt{\frac{x-b}{b-a}}\right).
\]
This follows from $T_k(1+u)=\cosh(k\arcosh(1+u))$ and
$\arcosh(1+u)\le\sqrt{2u}$. The inequalities
$x^2\ge4b(x-b)$ and $m\sqrt{b(b-a)}\ge k$ imply
$|e^{-mx}P(x)|\le\rho$.
Finally, $\arcosh((b+a)/(b-a))=2\artanh\sqrt{a/b}$, so
\[
 \rho_j\le2e^{-2k_j\sqrt{a_j/b}}\le2e^{-2q_j}
 =\delta_j^2/2.
\]
\end{proof}

Write
\[
 B_0=m(b-a),\quad v=\frac{a}{b-a},\quad h=\frac{k}{\sqrt{B_0}},
 \quad A=ma,\quad
 u_\ell=|c_\ell|\ell!/m^\ell\quad(0\le\ell\le k).
\]
The notation $B_0$ in this section is unrelated to the failure bound $B$.
The signs alternate: $c_\ell=(-1)^{\ell+1}|c_\ell|$ and $u_0=1$.

\begin{lemma}[Integral and range bounds]\label{lem:integral}
Let $I=\sum_{\ell=0}^k u_\ell$. Then $\|f\|_\infty\le I$ and the
range length of $f$ is at most $I$. Moreover,
\begin{align}
 I&\le1+2\sqrt\pi h e^{h^2},\label{eq:coarse}\\
 I&\le1+2\sqrt\pi h\exp(h^2-2k\sqrt v+A),\label{eq:refined}\\
 I&\le\frac1{1-\theta}\quad\text{if }\theta=\frac{k}{m\sqrt{ab}}<1.
 \label{eq:derivative}
\end{align}
\end{lemma}
\begin{proof}
Positivity of the coefficients of $T_k(z_0+x)$ for $z_0=(b+a)/(b-a)>1$
gives the exact factorial-weight identity
\[
 I=\int_0^\infty e^{-s}\frac{T_k(z_0+2s/B_0)}{T_k(z_0)}\,ds.
\]
For $t,u\ge0$, the identity $\arcosh(1+2t)=2\arsinh\sqrt t$ gives
\[
 \arcosh(1+2t+2u)-\arcosh(1+2t)
 \le2(\sqrt{t+u}-\sqrt t)\le2\sqrt u.
\]
Also $\cosh(x+y)/\cosh x\le e^y$ for $x,y\ge0$.
Thus the integral is at most
\[
 \int_0^\infty e^{-s+2h\sqrt s}\,ds
 \le1+2\sqrt\pi h e^{h^2},
\]
and, retaining the subtraction, at most
\[
 \int_0^\infty e^{-s+2h(\sqrt{s+A}-\sqrt A)}\,ds
 =1+\sqrt\pi h e^{(h-\sqrt A)^2}\operatorname{erfc}(\sqrt A-h).
\]
Using $\operatorname{erfc}\le2$ proves \eqref{eq:refined} because
$h\sqrt A=k\sqrt v$.
Alternatively, concavity of $\arcosh$ bounds the log of the polynomial
ratio by $\theta s$, proving \eqref{eq:derivative}.
For the range claim, use the convention that a maximum over an empty
set of nonnegative corrections equals zero. The largest positive correction is at most
$\max_{\ell\text{ odd}}u_\ell$, and the largest negative correction
is at most $\max_{\ell\text{ even},\ell\ge2}u_\ell$.
If the latter is at most one, the range length is at most
$1+\max_{\ell\text{ odd}}u_\ell\le I$; otherwise it is at most the sum
of these two maxima, also at most $I$. Include the values $f(0)=0$ and
$f(\ell)=1$ for $\ell>k$ in this argument.
\end{proof}

\begin{lemma}[Individual coefficients]\label{lem:coefficient}
Put $t=k^2/B_0$ and $w=k^2v$. For $1\le\ell\le k$,
\begin{equation}\label{eq:individual}
 u_\ell\le\rho\frac{(4t)^\ell\ell!}{(2\ell)!}
                   \exp\left(\frac{2w}{2\ell+1}\right),\qquad
 \frac{(4t)^\ell\ell!}{(2\ell)!}\le e^t.
\end{equation}
In particular, at scale zero,
\begin{equation}\label{eq:coefficientF}
 \|f_0\|_\infty\le\mathcal F(t,w):=
 1+\frac1{32}\max\left\{e^{t+2w/17},
 \max_{1\le\ell\le7}\frac{(4t)^\ell\ell!}{(2\ell)!}e^{2w/(2\ell+1)}\right\}.
\end{equation}
\end{lemma}
\begin{proof}
The polynomial identity
\[
 T_k(1+2x)=1+\sum_{\ell=1}^k
 \frac{4^\ell}{(2\ell)!}k^2\prod_{r=1}^{\ell-1}(k^2-r^2)x^\ell
\]
follows, for example, by substitution into
$(z^2-1)T_k''+zT_k'-k^2T_k=0$ and $T_k(1)=1$.
More explicitly, if $T_k(1+2x)=\sum_{\ell\ge0}a_\ell x^\ell$,
the differential equation becomes
$x(1+x)y''+(x+1/2)y'-k^2y=0$ and gives
$a_{\ell+1}=2(k^2-\ell^2)a_\ell/((\ell+1)(2\ell+1))$.
Starting with $a_0=1$ proves the displayed product formula, including its
termination at degree $k$.
Its coefficients are bounded coefficientwise by those of
$\cosh(2k\sqrt x)$. Substituting $x=v+s/B_0$ preserves this order.
The latter entire power series has nonnegative coefficients, so its
translated coefficients are obtained by absolutely convergent sums.
The coefficient of $s^\ell$, multiplied by $\ell!$, is at most
\[
 \frac{(4t)^\ell\ell!}{(2\ell)!}
 \sum_{r\ge0}\frac{(4w)^r}{r!}
 \frac{(\ell+r)!(2\ell)!}{\ell!(2\ell+2r)!}
 \le\frac{(4t)^\ell\ell!}{(2\ell)!}e^{2w/(2\ell+1)},
\]
since each factor of the last factorial ratio is
$1/[2(2\ell+2s-1)]\le1/[2(2\ell+1)]$.
Multiply by $\rho$ to obtain the first assertion.
For the second, midpoint concavity of $\log$ gives
\[
 \log\prod_{i=1}^\ell(i-1/2)\ge\int_0^\ell\log x\,dx
 =\ell\log\ell-\ell.
\]
Because this product is $(2\ell)!/(4^\ell\ell!)$, the desired quantity
is at most $(et/\ell)^\ell\le e^t$, using $1+\log y\le y$.
Finally $\rho_0\le1/32$ and $|f(\ell)|\le1+u_\ell$.
Split the indices at $\ell=8$ to obtain \eqref{eq:coefficientF}.
\end{proof}

\section{From moments to distribution-uniform error bounds}
Suppose $F_j$ bounds $\|f_j\|_\infty$, and $V_j$ bounds
$\Var(f_j(N))$ for every Poisson mean. Lemma~\ref{lem:integral} permits
$V_j=I_j^2/4$, while the coefficient bound permits $V_0\le F_0^2$.
Indeed, a variable in an interval of length $d$ has variance at most $d^2/4$:
take its second moment about the interval midpoint and minimize over centers.
Independence gives
\begin{equation}\label{eq:variance}
 \Var Z_j\le S V_j\quad(S<\infty),\qquad
 \Var Z_j\le F_j^2\E D\quad\text{in general}.
\end{equation}
For a countable alphabet, truncate it. An omitted mass $\tau$ contributes
in absolute value at most $F_j M_{\rm tail}$, where
$M_{\rm tail}\sim\Poi(m\tau)$. Its second moment is at most
$F_j^2(m\tau+(m\tau)^2)\to0$, justifying the limit and all moment formulas.
Also $\sum_i\E|f_j(N_i)|\le F_j\sum_i\Prb(N_i>0)\le F_jm$,
so the expectation series is absolutely summable. The expression
$\sum_i\min(p_i/a_j,1)$ is finite, being at most $1/a_j$.

We use the one-sided variance bound
\begin{equation}\label{eq:cantelli}
 \Prb(X-\E X\ge t)\le\frac{\Var X}{\Var X+t^2}\quad(t>0),
\end{equation}
and the same bound for the lower tail. To prove it, apply Markov to
$(X-\E X+a)^2$ and minimize $(\Var X+a^2)/(t+a)^2$ at
$a=\Var X/t$.
When the variance is zero the statement follows directly. This argument
applies to either weak one-sided tail, so equality at the decision threshold
introduces no exceptional event.

Define the dimensionless quantities
\[
 v_j=\frac{V_j}{n\delta_j^2},\quad w_j=\frac{F_j^2}{n\delta_j^2},\quad
 d_j=2-\delta_j/2,\quad \kappa_j=14-\delta_j/2-8\delta_j^2\ge13.375.
\]
If $S\le n$, the upper error $\widetilde S>S$ can only arise at a scale.
Its threshold is $(1+\delta_j)S+\delta_jn$, whereas
$\E Z_j\le(1+\rho_j)S$. Writing $u=S/n\in(0,1]$, we have
\[
 \frac{u}{(1+(1-\delta_j/2)u)^2}\le\frac1{(2-\delta_j/2)^2}.
\]
Indeed, the derivative of $u/(1+cu)^2$ is $(1-cu)/(1+cu)^3\ge0$ for
$0\le u\le1$ and $c=1-\delta_j/2\in[7/8,1)$.
Apply \eqref{eq:cantelli}, \eqref{eq:variance}, and a union bound:
\begin{equation}\label{eq:upperbound}
 \Prb(\widetilde S>S)\le U:=\sum_{j=0}^J\frac{v_j}{v_j+d_j^2}.
\end{equation}

For power, fix $1\le r\le n$ with $\Delta_r\ge\epsilon$.
If $\mu=\E D>2n$, then $\Var D\le\mu$ and
\begin{equation}\label{eq:largeocc}
 \Prb(\widetilde S\le r)\le\Prb(D\le n)
 \le\frac\mu{(\mu-n)^2}\le\frac2n.
\end{equation}
Otherwise $\mu\le2n$. There are three exhaustive cases.
\begin{enumerate}
\item If $p_{(r)}\ge b/2$, observing all of the first $r$ categories and
any tail category implies $D>r$. A union bound gives
\begin{equation}\label{eq:heavy}
 \Prb(D\le r)\le r e^{-mb/2}+e^{-m\Delta_r}
 \le H(L):=e^{-31L-1}+e^{-C}.
\end{equation}
Here $n/L\ge1$, since $n\ge1+\log n$.
\item If $p_{(r)}<a_0$, the tail alone gives
$\E Z_0\ge(1-1/32)\epsilon/a_0\ge7.75n$.
The rejection threshold at $r$ is at most $1.5n$, so its margin is at least
$6.25n=25\delta_0n>\kappa_0\delta_0n$.
\item If $a_0\le p_{(r)}<b/2$, choose $j$ with
$a_j\le p_{(r)}<2a_j$. Such a $j\le J$ exists. For every tail mass $x$,
$\min(x/a_j,1)\ge x/p_{(r)}$, while each of the first $r$ masses
contributes one. Since $a_j\delta_j=\epsilon/(32n)$,
\[
 \E Z_j\ge(1-\rho_j)(r+16\delta_j n).
\]
Subtracting $(1+\delta_j)r+\delta_jn$ leaves a margin at least
\[
 \bigl[14\delta_j-\rho_j(1+16\delta_j)\bigr]n
 \ge\kappa_j\delta_jn.
\]
\end{enumerate}
These are cases for one fixed distribution, so their error bounds are
maximized, not summed. Consequently
\begin{align}
 \Prb(\widetilde S\le r)&\le R:=
 \max\left\{H(L),\frac2n,\max_j\frac{2w_j}{2w_j+\kappa_j^2}\right\},
 \label{eq:powergeneral}\\
 S\le n\quad\Longrightarrow\quad
 \Prb(\widetilde S\le r)&\le R_f:=
 \max\left\{H(L),\max_j\frac{v_j}{v_j+\kappa_j^2}\right\}.
 \label{eq:powerfinite}
\end{align}
For the second statement, $\mu>2n$ is impossible and
$\Var Z_j\le nV_j$.

To prove the joint guarantee, let $t=\min(\eff_\epsilon(p),n)$.
If $t=1$, its lower constraint is automatic. Otherwise take $r=t-1$;
then $\Delta_r>\epsilon$ and $\widetilde S>r$ implies $\widehat S\ge t$.
In this case $1\le r\le n-1$, so the previously proved power inequality
applies. Since $S\ge1$, clipping the output below at one preserves its
upper bound whenever $\widetilde S\le S$.
If $S\ge n$, the upper constraint $\widehat S\le S$ is automatic.
If $S<n$, combine \eqref{eq:upperbound} and \eqref{eq:powerfinite}.
Thus a common error bound for both binary promises and joint coverage is
\begin{equation}\label{eq:common}
 B(L)=\max\{U+R_f,R\}.
\end{equation}
This use of $R_f$ is essential: a binary-test guarantee alone does not imply
the same joint-coverage guarantee.

\section{Uniform parameter certification}
The following covers all $n$ and $\epsilon$, not a grid of sampled inputs.
The only remaining numerical obligations are finitely many exact rational
inequalities. All intervals between endpoints and the infinite tail are
covered by the displayed analytic inequalities. Appendix~\ref{app:arithmetic}
proves the soundness of the rational enclosures used for these obligations.

\subsection{Well-defined parameters and the counting branch}
For $C\in\{144(\log2)^2,70\}$ and $L\ge15$,
\[
 69<C\le70,\qquad \frac b{2a_0}=\frac{256L^2}{C}>1,\qquad
 a_j\le b/2,\qquad b\le\frac{64L^2}{Ce^{L-1}}<1.
\]
The last function decreases for $L>2$. Moreover,
\[
 k_j\le\alpha_j L+1,\quad
 \alpha_j=\sqrt{512/C}\,\frac{q_j}{2^{j/2}},\quad
 k_j<6L<\frac{64L}{\sqrt2}\le m\sqrt{b(b-a_j)}.
\]
Here $(j+3)2^{-j/2}$ decreases for integers $j\ge0$.
For explicit elementary bounds, $693/1000<\log2<694/1000$ implies
$69<144(\log2)^2<70$, while
$\sqrt{512/69}<11/4$ gives
$k_j/L<(11/4)\cdot3\cdot(7/10)+1/15<6$.
At $L=15$, the bound $b<1$ already follows from $e^{14}>1000$.
Thus Lemma~\ref{lem:expect} applies throughout.

If $\Delta_r\ge\epsilon$, the discovery process has unseen mass at least
$\epsilon$ before its $(r+1)$st distinct label. Its inter-discovery times
can be coupled as $E_k/q_k\le E_k/\epsilon$, where the $E_k$ are independent
unit exponentials. This follows by the memoryless property of a marked
Poisson process; the next unseen label is sampled independently of that
waiting time from normalized unseen masses. Hence
\[
 \Prb(D\le r)\le\Prb\{\Poi(m\epsilon)\le r\}
 \le\Prb\{\Poi(Cn/L)\le n\}.
\]
For $x>1$, exponential Markov gives
$\Prb\{\Poi(nx)\le n\}\le e^{-ng(x)}$, $g(x)=x-1-\log x$.
When $n<100$, $L<6$; when $n\ge100$ and $L<L_0\le20$, $L<20$.
Therefore, for every listed pair, counting failure is at most
\begin{equation}\label{eq:countcert}
 \max\{e^{-g(C/6)},e^{-100g(C/L_0)}\}<10^{-3}.
\end{equation}
The upper error in this branch is zero because $D\le S$.
No numerical logarithm is needed for \eqref{eq:countcert}:
$e^5>100$ implies $1+\log100<6$;
$e^{5/2}>23/2$ implies $g(C/6)>g(23/2)>8$;
and $e^{29/20}>69/20$ implies $g(C/L_0)>g(69/20)>1$.
Finally $e^8>1000$. These four inequalities follow from the positive Taylor
partial sums of degrees $8,8,5,12$, respectively, and are verified as
exact fractions by the certificate.

\subsection{Complete finite intervals: $15\le L\le32$}
Partition this interval into the 170 closed intervals
$[a,b]=[i/10,(i+1)/10]$, $i=150,\ldots,319$.
To avoid confusing $b$ here with the polynomial endpoint, use the following
slab-specific formulas solely in this subsection. Include all integer scales
$j\ge0$ with $2^j\le256b^2/C$. Even if a scale does not exist everywhere in
the slab, including its nonnegative error bound is safe. Define
\begin{align*}
 \bar k_j&=\alpha_j b+1,&
 \bar u_j&=\min\left\{\tfrac12,\frac{C2^j}{512a^2}\right\},&
 T_j&=\frac{\bar k_j^2}{64a(1-\bar u_j)},\\
 W_j&=\frac{\bar k_j^2\bar u_j}{1-\bar u_j},&
 \Theta_j&=\frac{\bar k_j}{\sqrt{8C2^j}},&
 A_j^+&=\frac{C2^j}{8a}.
\end{align*}
Throughout the slab, $h_j^2\le T_j$, $k_j^2a_j/(b_{\rm poly}-a_j)\le W_j$,
$ma_j\le A_j^+$, and $k_j\sqrt{a_j/(b_{\rm poly}-a_j)}\ge q_j$.
By Lemma~\ref{lem:integral} an upper bound for $I_j$ is
\[
 \bar I_j=\min\left\{
 1+2\sqrt{\pi T_j}e^{T_j},\quad
 1+2\sqrt{\pi T_j}e^{T_j-2q_j+A_j^+},\quad
 \frac1{1-\Theta_j}\ \text{if }\Theta_j<1\right\}.
\]
Use $\bar F_j=\bar I_j$, $\bar V_j=\bar I_j^2/4$ for $j\ge1$;
at scale zero use
\[
 \bar F_0=\min\{\bar I_0,\mathcal F(T_0,W_0)\},\qquad
 \bar V_0=\min\{\bar I_0^2/4,\mathcal F(T_0,W_0)^2\}.
\]
Insert $v_j\le\bar V_j e^{1-a}\delta_j^{-2}$,
$w_j\le\bar F_j^2 e^{1-a}\delta_j^{-2}$, $H(L)\le H(a)$, and
$2/n\le2e^{1-a}$ into \eqref{eq:common}.
These formulas are monotone in their error arguments. They cover all
degree-ceiling and scale-floor discontinuities without evaluating those
discontinuities numerically.

\subsection{The entire half-line: $L\ge32$}
Put $H=32$. For $0\le j\le7$, define
\[
 \xi_j=\frac{C2^j}{512},\qquad
 t_j=\frac{\alpha_j^2}{64}=\frac{8q_j^2}{C2^j},\qquad
 d_j^+=\frac{(\alpha_jH+1)^2}{64H(1-\xi_j/H^2)}-t_jH.
\]
Then, for every $L\ge H$,
\begin{equation}\label{eq:affine}
 h_j^2\le\frac{(\alpha_jL+1)^2}{64L(1-\xi_j/L^2)}
 \le t_jL+d_j^+.
\end{equation}
Indeed, after subtracting $t_jL$, the middle expression equals
\[
 \frac{2\alpha_j+(1+\alpha_j^2\xi_j)/L}
      {64(1-\xi_j/L^2)},
\]
which decreases with $L$. At scale zero also
\[
 k_0^2\frac{a_0}{b-a_0}\le
 W_0^+=\frac{(\alpha_0+1/H)^2 C}{512(1-C/(512H^2))}.
\]
Now the particular constant emerges:
\begin{equation}\label{eq:critical}
 2t_0=\frac{144(\log2)^2}{C}\le1.
\end{equation}
For $C=144(\log2)^2$ this is exact equality; it is not a numerical
comparison of nearly equal floating-point numbers.
Using \eqref{eq:coefficientF}, the bound
$e^{1-L}\mathcal F(t_0L+d_0^+,W_0^+)^2$ is nonincreasing on $L\ge H$.
For its exponential term this follows from $2t_0\le1$; for the seven
polynomial terms it follows from $H\ge2\cdot7$.
To see this without an asymptotic argument, write each candidate as
$1+A e^{t_0L}$ or $1+A(t_0L+d_0^+)^\ell$, $A>0$.
For the first, the square times $e^{-L}$ expands into three exponentials
with nonpositive exponents. For the second, its logarithmic derivative
after multiplication by $e^{-L}$ is at most $2\ell/L-1\le0$.
The maximum of finitely many nonincreasing functions is nonincreasing.
Thus its value at $H$ bounds the entire scale-zero contribution.

For $1\le j\le7$, let
\[
 z_j=\min\left\{d_j^+,\ d_j^+-2q_j+\frac{C2^j}{8H}\right\}.
\]
Then $I_j\le1+2\sqrt{\pi(t_jL+d_j^+)}e^{t_jL+z_j}$.
The square of this bound multiplied by $e^{1-L}$ is nonincreasing on
$L\ge H$: expand its three terms and use
$H(1-2t_j)>1$, verified by the certificate. Hence each can be bounded at $H$.
For the square-root cross term its logarithmic derivative is at most
$1/(2L)-(1-t_j)<0$; for the last term it is at most
$1/L-(1-2t_j)<0$. The constant term times $e^{-L}$ decreases as well.
Here $d_j^+>0$ follows from the explicit expression below
\eqref{eq:affine}.

Finally, for every active $j\ge8$,
\[
 h_j^2\le\beta^2L,\qquad
 \beta^2=\frac{(\alpha_8+1/H)^2}{32},\qquad
 I_j\le1+2\sqrt\pi\beta\sqrt L e^{\beta^2L}.
\]
This uses $m(b-a_j)\ge32L$ and the monotonicity of $\alpha_j$.
Since $\delta_J^{-1}\le1024L^2/C$,
\begin{equation}\label{eq:tail}
 \sum_{j=8}^J\frac{F_j^2}{n\delta_j^2}
 \le T(L):=\frac43\left(\frac{1024L^2}{C}\right)^2
 e^{1-L}\left(1+2\sqrt\pi\beta\sqrt L e^{\beta^2L}\right)^2
 \le T(H).
\end{equation}
The derivative signs follow by expanding into powers $L^4,L^{9/2},L^5$
and using $H(1-2\beta^2)>5$.
The geometric-sum factor is justified by
$\sum_{j=8}^J\delta_j^{-2}\le\sum_{j=0}^J16\cdot4^j
\le(4/3)\delta_J^{-2}$, with the empty sum interpreted as zero.
For the low scales conservatively take $V_j\le F_j^2$.
In \eqref{eq:upperbound}, the omitted scales contribute at most $T(H)$;
in \eqref{eq:powerfinite} at most $T(H)/81$; in
\eqref{eq:powergeneral} at most $2T(H)/81$, since $\kappa_j>9$.
For the upper-error contribution we also used $d_j\ge15/8>1$.
This yields the certified half-line bounds
\[
 \sup_{L\ge32}B(L)<
 \begin{cases}
 0.070579,&C=144(\log2)^2,\\
 0.051553,&C=70.
 \end{cases}
\]

\subsection{Certificate and proof status}
The verifier \texttt{exact\_certificate.py}, together with the small arithmetic
module \texttt{rational\_interval.py}, evaluates the displayed expressions
using integer endpoints with denominator $2^{160}$. Every division and
square root is rounded outward by integer floor/ceiling operations;
exponentials, $\log2$, and $\pi$ use rational series with proved remainders.
Appendix~\ref{app:arithmetic} proves these enclosure rules. No numerical
library, Monte Carlo sample, or floating-point transcendental function is
used in these certificates.

For each $C$, ten slab bounds are grouped into each row of
Table~\ref{tab:slabs}. The entries are rational numbers rounded upward,
not estimates. The file \texttt{exact\_certificate.json} contains the rational
inputs and output inequalities. A separate Fraction-based audit
recalculates the risk aggregation from those inputs and checks the printed
claims. The maxima over the applicable rows, together with the counting
bound $1/1000$, give Table~\ref{tab:constants}. This proves
Theorem~\ref{thm:main}.

\begin{table}[ht]
\centering
\begin{tabular}{lrr}
\toprule
Entire $L$ interval & $C=144(\log2)^2$ & $C=70$\\
\midrule
$[15,16]$ & 0.392089 & 0.376264 \\
$[16,17]$ & 0.317480 & 0.282612 \\
$[17,18]$ & 0.247691 & 0.232405 \\
$[18,19]$ & 0.147401 & 0.131984 \\
$[19,20]$ & 0.105619 & 0.090185 \\
$[20,21]$ & 0.087747 & 0.072224 \\
$[21,22]$ & 0.079410 & 0.063748 \\
$[22,23]$ & 0.074918 & 0.059079 \\
$[23,24]$ & 0.072069 & 0.056023 \\
$[24,25]$ & 0.070445 & 0.054169 \\
$[25,26]$ & 0.068530 & 0.052005 \\
$[26,27]$ & 0.067068 & 0.050277 \\
$[27,28]$ & 0.065872 & 0.048801 \\
$[28,29]$ & 0.064857 & 0.047495 \\
$[29,30]$ & 0.063978 & 0.046315 \\
$[30,31]$ & 0.063209 & 0.045238 \\
$[31,32]$ & 0.062534 & 0.044245 \\
$[32,\infty)$ & 0.070579 & 0.051553 \\
\bottomrule
\end{tabular}
\caption{Exact rational upper bounds on $B(L)$ over complete intervals.
The last row is obtained by analytic monotonicity on an unbounded interval.}
\label{tab:slabs}
\end{table}

This is a mathematical proof with a finite exact-arithmetic component.
It is not an end-to-end formalization in a proof-assistant kernel. Its
analytic lemmas and the source of the certificate verifier remain explicit
objects for mathematical review; a successful program run alone is not
offered as a substitute for those lemmas.

\section{Confidence amplification and its actual cost}
\begin{theorem}[Confidence amplification]\label{thm:amplification}
Let $1/2<q\le1$ be a uniform lower bound on single-run success, independent of
$p,n,\epsilon$. Run an odd number $r$ of independent copies, each with its
own Poisson sample count and fresh independent data.
Use majority voting for the Yes/No output and the median for
$\widehat S$. Either aggregated failure is at most
\begin{equation}\label{eq:binomial}
 \sum_{k=(r+1)/2}^r {r\choose k}(1-q)^kq^{r-k}
 \le [4q(1-q)]^{r/2}
 \le e^{-2r(q-1/2)^2}.
\end{equation}
\end{theorem}
\begin{proof}
The case $q=1$ is immediate, so assume $q<1$.
If more than half of the support-bound outputs belong to the fixed interval
$[\min(\eff_\epsilon(p),n),S]$, their median belongs to that interval.
This interval is nonempty: for finite $S$, $\eff_\epsilon(p)\le S$,
and for infinite $S$ there is no finite upper constraint.
The corresponding implication for binary majority voting is immediate.
If $X_i$ is a failure indicator with probability at most
$1-q$, exponential Markov at $\lambda=\log(q/(1-q))$ bounds
$\Prb(\sum X_i\ge r/2)$ by
\[
 e^{-\lambda r/2}[q+(1-q)e^\lambda]^r=[4q(1-q)]^{r/2}.
\]
The last inequality in \eqref{eq:binomial} uses
$4q(1-q)=1-4(q-1/2)^2$ and $1-x\le e^{-x}$.
The exact binomial bound follows by coupling each failure indicator to an
independent Bernoulli variable of parameter $1-q$.
\end{proof}

Expected total samples equal $rm=rCn/(\epsilon L)$.
For target failure $\eta$, the sufficient repetition count is the least odd
integer at least $\log(1/\eta)/(2(q-1/2)^2)$; exact binomial tails improve it.
For example, exact rational calculations give
\begin{center}
\begin{tabular}{rrr}
\toprule
guaranteed $q$ & repeats $r$ & failure upper bound\\
\midrule
$2/3$&47&0.009002915\\
$3/5$&133&0.009785038\\
$0.92$&5&0.0045252608\\
\bottomrule
\end{tabular}
\end{center}
Therefore $(C,L_0)=(70,20)$ with five copies has expected coefficient
$5C=350$, while relying only on $q\ge2/3$ and 47 copies would give
coefficient $3290$. Minimizing the single-run coefficient alone is not the
same optimization problem as minimizing total samples at a fixed confidence.
The two amplified guarantees are individual statements. If one requires
the binary output and the support bound to be correct simultaneously, a
union bound gives twice the displayed failure bound; allocate $\eta/2$ to
each when needed.

Merely having success $>1/2$ for each separate input is insufficient for a
uniform constant repetition budget. One needs a fixed margin
$\gamma=q-1/2>0$. Reusing the same data with fresh random seeds does not
provide the independent-copy hypothesis. Bootstrap resampling likewise does
not supply this guarantee about the unknown underlying distribution.

\begin{corollary}[Deterministic sample cap]\label{cor:cap}
In a single run, draw $M\sim\Poi(m)$ before accessing the sample oracle.
If $M>K:=\lceil2m\rceil$, return any fixed allowed outputs without querying
the oracle; otherwise perform the original run. This algorithm uses at most
$K$ samples, has expected sample use at most $m$, and has each error bound
of Theorem~\ref{thm:main} increased by less than $10^{-9}$.
All convenient success guarantees in Table~\ref{tab:constants} remain valid.
In particular, five independent capped runs with $(C,L_0)=(70,20)$ have
individual binary and support-bound failure at most $0.0045252608$,
joint failure of either output at most $0.0090505216$, expected total samples
at most $350n/(\epsilon L)$, and total samples always at most
$700n/(\epsilon L)+5$.
\end{corollary}
\begin{proof}
Couple the capped run to the original run. Their outputs agree unless
$M>K$. The Poisson moment generating function and exponential Markov at
$\log2$ give
\[
 \Prb(M>K)\le\Prb(M\ge2m)
 \le\exp\{-m(2\log2-1)\}.
\]
Since $n/L\ge1$ and $\epsilon<1$, $m>C\ge C_*>69$.
The elementary bounds above give
$m(2\log2-1)>69(2\cdot693/1000-1)>24$ and
$e^{24}=(e^8)^3>10^9$. A union bound proves the error increase.
The rational slack in every row of Table~\ref{tab:constants} exceeds
$10^{-9}$ relative to its stated convenient success guarantee.
For the five-run claim use $q=23/25$ in \eqref{eq:binomial}:
the exact failure is at most $44192/9765625=0.0045252608$ per task.
Apply one union bound for both outputs. Finally,
$5\lceil2m\rceil\le10m+5$ and the expected number queried is at most $5m$.
No independence claim is made conditional on the no-abort event; the proof
uses the unconditional coupling instead. Different capped runs use
independent randomness and samples, as required by amplification.
\end{proof}

\section{What is optimized, and what remains open}
The new constant is over fourteen times smaller than $1025$, and the
polynomial branch starts near $n=e^{15}\approx3.27\times10^6$ for $L_0=16$,
instead of near $e^{1023}$. For $L_0=15$, the threshold is $e^{14}$;
for $L_0=20$, it is $e^{19}$. These are thresholds in $n$, not guarantees
that the required sample budget is practical for every $\epsilon$.

Equation~\eqref{eq:critical} explains why this particular proof gives
$144(\log2)^2$. Below it, the scale-zero leading upper bound used here grows
like $\exp((144(\log2)^2/C-1)L)$. This is a limitation of this bound, not a
counterexample to the algorithm. A sharper coefficient or distribution-aware
variance argument might lower $C$ further. Nor have we proved that $L_0=16$
is the smallest threshold achieving $2/3$ success. Parameters have a tradeoff;
there is no justified claim here that one pair minimizes both globally.

A weak necessary condition illustrates the remaining gap. If $L_0>1$, the
input $n=1$ uses the counting branch. On the uniform distribution on $K$
labels take $\epsilon=1-1/K$, so $\Delta_1=\epsilon$. As $K\to\infty$,
$m\to C$, collisions vanish, and the error probability tends to
$e^{-C}(1+C)$. Hence uniform success at least $2/3$ requires
$e^{-C}(1+C)\le1/3$, or $C\ge2.2892814\ldots$.
A uniform success constant strictly above $1/2$ requires
$C>1.6783469\ldots$. These necessary conditions are far below our sufficient
constant and are not tight optimality results.

\paragraph{Implementation scope.}
The proof uses the real mean in \eqref{eq:model}, exact polynomial weights,
and independent Poissonized samples. It does not certify arbitrary
floating-point approximations, an integer-rounded mean with recomputed
weights, or old experiments with other constants. If only integer Poisson
means are supported, draw $\Poi(\lceil m\rceil)$ samples and independently
retain each with probability $m/\lceil m\rceil$; the retained counts have
exactly the law in \eqref{eq:model}. Compute weights from $m$ and use retained
counts. This costs less than one extra expected sample per run.
These mathematical guarantees do not by themselves establish numerical
stability of a floating-point implementation or measured performance on
empirical distributions.

\begin{thebibliography}{2}
\bibitem{wang} Qisheng Wang. \emph{Optimal Support Size Testing}.
arXiv:2610.18723v1, 2026. \url{https://arxiv.org/abs/2610.18723v1}.
\bibitem{dlmf} NIST Digital Library of Mathematical Functions.
Section 18.5, Explicit Representations.
\url{https://dlmf.nist.gov/18.5}. Accessed October 21, 2026.
\end{thebibliography}
\appendix
\section{Exact arithmetic certificate and its soundness}\label{app:arithmetic}

This appendix specifies all numerical operations needed in the proof and
proves that their outputs enclose the corresponding real expressions.
The numerical part of the argument is a finite sequence of integer
inequalities; it does not assume correctness of a floating-point error
estimate. The source files and complete rational records accompany the paper.

\subsection{Dyadic representation and elementary operations}
Fix $P=160$ and $D=2^P$. An enclosure is a pair of integers $(a,b)$,
$a\le b$, representing $[a/D,b/D]$.
An exact fraction $s/t$, $t>0$, is enclosed by
\[
 \left(\left\lfloor\frac{sD}{t}\right\rfloor,
       \left\lceil\frac{sD}{t}\right\rceil\right).
\]
Addition and negation act exactly on endpoints. To multiply two enclosures,
take the minimum and maximum of the four integer endpoint products, divide
by $D$, and round the lower endpoint down and the upper endpoint up.
If $0<a\le b$, the reciprocal enclosure is
\[
 \left(\left\lfloor\frac{D^2}{b}\right\rfloor,
       \left\lceil\frac{D^2}{a}\right\rceil\right).
\]
For a strictly negative interval use sign reversal; division is rejected if
the denominator interval contains zero. Integer powers are evaluated by
repeated squaring, using these same enclosure operations.
For square roots of $0\le a\le b$, let
\[
 u=\lfloor\sqrt{aD}\rfloor,\qquad v=\lfloor\sqrt{bD}\rfloor.
\]
The integer witnesses $u^2\le aD<(u+1)^2$ and
$v^2\le bD<(v+1)^2$ are checked explicitly. The resulting enclosure is
$(u,v)$ if $v^2=bD$, and $(u,v+1)$ otherwise.

For intervals $X_i=[a_i,b_i]$, the minimum and maximum of their unknown
values are enclosed, respectively, by
$[\min_i a_i,\min_i b_i]$ and $[\max_i a_i,\max_i b_i]$.
These rules do not require independence among the unknown values.
Likewise, reusing an uncertain expression may widen an interval but cannot
invalidate its enclosure. A strict comparison $X<Y$ is accepted only when
the upper endpoint of $X$ is below the lower endpoint of $Y$.

\begin{lemma}[Elementary enclosure invariant]\label{lem:invariant}
Every elementary operation just specified encloses its real result whenever
its operands are valid enclosures and its stated domain conditions hold.
\end{lemma}
\begin{proof}
For multiplication, a bilinear function attains its extrema on a rectangle
at corners. Reciprocal is decreasing on each of its sign domains, square
root is increasing on $[0,\infty)$, and minimum and maximum are coordinatewise
increasing. The floor and ceiling definitions preserve the appropriate
one-sided bounds. Addition and negation are immediate. Induction on the
number of multiplications proves the integer-power claim.
\end{proof}

\subsection{Exponential enclosures}
For a nonnegative exact rational $x$, choose an integer $r\ge0$ so that
$y=x/2^r\le1/8$. With $N=40$, define the exact fractions
\[
 S_N(y)=\sum_{k=0}^N\frac{y^k}{k!},\qquad
 R_N(y)=\frac{y^{N+1}}{(N+1)!}\frac1{1-y/(N+2)}.
\]
Then
\begin{equation}\label{eq:exp-enclosure}
 S_N(y)\le e^y\le S_N(y)+R_N(y).
\end{equation}
Indeed, all terms are nonnegative, and the ratio of consecutive terms after
the first omitted one is at most $y/(N+2)<1$, so a geometric series bounds
the tail. Enclose both endpoints as dyadic fractions and square the result
$r$ times. This encloses $e^x$ by Lemma~\ref{lem:invariant}. For negative
$x$, enclose $e^{-x}$ first and take the reciprocal. For an input interval
$[a,b]$, enclose $e^a$ and $e^b$ separately and retain the outer endpoints,
using monotonicity of $\exp$.

All Taylor coefficients and remainders in the implementation are computed
with exact fractions. Thus the finite series length affects enclosure width,
not validity. Very small positive exponentials may be enclosed by an interval
whose lower endpoint is zero; this is a valid enclosure, not a replacement
of a positive real number by zero. A subsequent division by such an interval
would be rejected.

\subsection{Logarithm and circle-constant enclosures}
Let $z=1/3$ and $N=80$. Integrating the geometric series for $1/(1-z^2)$
gives
\[
 \log2=2\sum_{k=0}^\infty\frac{z^{2k+1}}{2k+1}.
\]
Its remainder after $N$ terms is positive and at most
\begin{equation}\label{eq:log-enclosure}
 \frac{2z^{2N+1}}{(2N+1)(1-z^2)}.
\end{equation}
Consequently this partial sum and remainder give an exact rational enclosure
for $\log2$. Squaring and multiplying by 144 gives one for $C_*$.
At the critical scale the identity $2t_0=1$ is simplified symbolically
before evaluation: for $C=C_*$,
\[
 t_j=\frac{(j+3)^2}{18\cdot2^j}.
\]
Thus the endpoint sign condition does not depend on rounding a number close
to one.

To enclose $\pi$ we use Machin's identity
\begin{equation}\label{eq:machin}
 \pi=16\arctan(1/5)-4\arctan(1/239).
\end{equation}
Here is a direct verification. If $\theta=\arctan(1/5)$ and
$\phi=\arctan(1/239)$, double-angle formulas give
$\tan(2\theta)=5/12$ and $\tan(4\theta)=120/119$; therefore
$\tan(4\theta-\phi)=1$. Also $0<4\theta-\phi<4/5<\pi/2$:
positivity follows from $\theta>\phi>0$, and
$\theta<1/5$ follows from the integral formula for arctangent.
The elementary bound $\pi/2>1$ follows from
$\pi/4=\int_0^1(1+x^2)^{-1}dx>1/2$.
Injectivity of tangent on this interval gives $4\theta-\phi=\pi/4$.

For $z\in\{1/5,1/239\}$, the alternating expansion
\[
 \arctan z=\sum_{k=0}^{N-1}(-1)^k\frac{z^{2k+1}}{2k+1}+R_N,
 \qquad 0\le R_N\le\frac{z^{2N+1}}{2N+1},\qquad N=50,
\]
follows by integrating the finite geometric expansion of $1/(1+x^2)$;
$N$ is even. Combining these rational enclosures in \eqref{eq:machin}
encloses $\pi$.

\begin{lemma}[Soundness of the certificate arithmetic]\label{lem:soundness}
Every enclosure produced by \texttt{rational\_interval.py} for an expression
used in the certificate contains the exact real value of that expression.
Every strict numerical inequality accepted by \texttt{exact\_certificate.py}
is therefore a true real inequality.
\end{lemma}
\begin{proof}
The leaf constants are exact rationals, the logarithm enclosure
\eqref{eq:log-enclosure}, and the circle-constant enclosure
\eqref{eq:machin}. They are sound by the preceding arguments.
Structural induction on the expression tree uses Lemma~\ref{lem:invariant}
for elementary operations and \eqref{eq:exp-enclosure} for exponentials.
Every division checks separation from zero, every square root checks its
nonnegative domain and integer witnesses, and conditional use of
$1/(1-\Theta_j)$ requires a certified strict inequality $\Theta_j<1$.
The last assertion follows from disjoint interval endpoints.
\end{proof}

\subsection{Finite specification of the obligations}
The verification comprises the following explicitly bounded operations.
\begin{enumerate}
\item Establish $693/1000<\log2<694/1000$, $69<C_*<70$, the parameter
inequalities, and the four positive Taylor inequalities in the counting proof.
\item For each of the two values $C\in\{C_*,70\}$ and each integer
$i\in\{150,\ldots,319\}$, evaluate the slab expressions on
$[i/10,(i+1)/10]$. Include scales until the first $j+1$ for which
$2^{j+1}$ is certified greater than $256((i+1)/10)^2/C$.
This can include an unnecessary scale if a comparison is inconclusive;
including additional nonnegative error terms is conservative.
\item For each $C$, evaluate the eight low-scale expressions at $H=32$,
check their derivative conditions, and evaluate the omitted-scale bound.
For $C=C_*$ use the exact rational formula for $t_j$ above.
\item For each row of Table~\ref{tab:constants}, take maxima over all slabs
starting at or above its $L_0$, the corresponding half-line bound, and
$1/1000$. Check all three required error bounds against the printed rational
$B$. Check that $B+10^{-9}$ is still below the failure threshold associated
with the table's convenient success guarantee.
\item Compute each amplification probability as an exact finite sum of
fractions, and compare it with $1/100$.
\end{enumerate}
All loops have explicit finite limits, except scale inclusion, which terminates
because its fixed finite comparison bound is eventually exceeded by $2^j$.
The exact intervals are retained in the certificate file. In particular,
there is no extrapolation from a sampled distribution, no test of only
selected integer $n$, and no truncation of the half-line $L\ge32$.

\subsection{Independent risk aggregation and reproducibility}
For every slab the certificate records rational upper bounds
$v_j^+,w_j^+,h^+,o^+$ for $v_j,w_j,H(L),2/n$.
The separate script \texttt{audit\_exact.py} reconstructs, using the standard
exact Fraction type,
\begin{align*}
 U^+&=\sum_j\frac{v_j^+}{v_j^++d_j^2},\\
 R_f^+&=\max\left\{h^+,\max_j\frac{v_j^+}{v_j^++\kappa_j^2}\right\},\\
 R^+&=\max\left\{h^+,o^+,\max_j\frac{2w_j^+}{2w_j^++\kappa_j^2}\right\},\\
 B^+&=\max\{U^++R_f^+,R^+\}.
\end{align*}
These are valid bounds because the displayed rational functions are
coordinatewise nondecreasing for nonnegative arguments.
The auditor similarly reconstructs the half-line aggregation, checks that
the complete set of slabs is present, verifies the source hashes, and checks
all claimed thresholds. It also tests the arithmetic operations against
independent Fraction calculations and checks exact Chebyshev coefficient
identities. Those tests complement the mathematical soundness proof;
they do not replace it.

From the artifact directory, the commands are
\begin{verbatim}
python3 exact_certificate.py
python3 audit_exact.py
tectonic submission_proof.tex
\end{verbatim}
Only Python's standard library is needed for the first two commands.
The checker rejects float inputs and does not call a logarithm, exponential,
or trigonometric routine from a numerical library. Its checks raise explicit
exceptions and remain active under Python's optimization flag.

\paragraph{Scope of formalization.}
The theorem, reduction to a finite family of rational inequalities, and
soundness of the enclosure rules are mathematical proofs supplied in this
document. The artifact executes those finite checks. There is no claim that
the complete probability and real-analysis development has been encoded in
Lean, Coq, Isabelle, or another foundational proof assistant, and no new
axiom asserting the desired statistical theorem is introduced. A full
proof-assistant formalization would be an additional deliverable.

\end{document}
